\documentclass[10pt,a4paper]{amsart}
\usepackage[margin=19mm]{geometry}
\usepackage{amsmath,amssymb,mathtools}
\usepackage[scr=rsfs,cal=euler]{mathalfa}
\usepackage{xfrac}
\usepackage[unicode,hidelinks,bookmarks=false]{hyperref}
\newcommand{\ie}{i.e.\ }
\newcommand{\cf}{cf.\ }
\newcommand{\resp}{respectively\ }
\newcommand{\Subsection}{\subsection}
\providecommand{\nobreakdash}{\nobreak\hbox{-}}
\setlength{\emergencystretch}{2em}
\multlinegap0pt
\usepackage{bm}
\usepackage{bbm}
\usepackage{dsfont}
\usepackage{xspace}
\usepackage{cancel}
\usepackage{mathtools}
\usepackage{amsthm}
\makeatletter
\@ifundefined{theorem}{\newtheorem{theorem}{Theorem}}{}
\@ifundefined{lemma}{\newtheorem{lemma}[theorem]{Lemma}}{}
\makeatother
\title[Schramm-Loewner evolution contains a topological Sierpi\'nsk]{Schramm-Loewner evolution contains a topological Sierpi\'nski carpet when $\kappa$ is close to 8}
\author{Haoyu Liu, Zijie Zhuang}
\date{Source: arXiv:2506.09609v2 (CC BY 4.0)}
\begin{document}
\maketitle

\noindent\emph{Source and licence.} Schramm-Loewner evolution contains a topological Sierpi\'nski carpet when $\kappa$ is close to 8. Version arXiv:2506.09609v2 (CC BY 4.0), distributed under the Creative Commons Attribution 4.0 International licence, as recorded in the arXiv licence metadata for this exact version and shown on the abstract page https://arxiv.org/abs/2506.09609v2. Full rights notice: licenses/arxiv-2506.09609-CC-BY-4.0.txt.\par\smallskip

\noindent\textit{Editorial scope.} Counted unit: the Lemma whose source label is \texttt{lem:length-bound} in the article (source0.tex lines 1035--1037, source label \texttt{lem:length-bound}), delivered together with the article's own complete original proof of it (source0.tex lines 1039--1068, one \texttt{proof} environment of 30 lines). No mathematics is written, repaired or re-derived: statement and proof are the article's own text line for line. Required context retained uncounted: condition:fractal-2 (lines 802--806). Every definition, display and label of those lines is retained unchanged. Omitted from this package and not redistributed: 1--801, 807--1034, 1038--1038, 1069--1273. Nothing inside a retained excerpt is omitted or altered: every retained body is delivered exactly as the article's lines are written, so no omission receipt of any kind accompanies this package. Citations: the retained excerpts carry no citation command at all, so no bibliography entry of the article is transcribed and no reference is redistributed. Reference adaptation: none was needed and none was made; every reference inside the delivered unit resolves inside it, so no unresolved reference or citation is produced. Printed slips kept literal: no printed word, symbol, hypothesis or number of the counted statement or of its proof was altered. Layout and class: the article's own class is \texttt{article} with packages including xy, inputenc, graphics, graphicx, amsfonts, amssymb, amsmath, xcolor, mathrsfs, amstext, amsbsy, amsopn, amscd, amsxtra; the wrapper is the lane's pinned \texttt{amsart} editorial wrapper at 10pt on A4, which restates the theorem-like environments the retained text needs and adds the article's own macro definitions that the retained text needs (none), with extra wrapper packages (bm, bbm, dsfont, xspace, cancel, mathtools). The delivered numbering is the unit's own. Assets: no retained span contains a figure environment, a graphics-inclusion command, a PDF-page inclusion, a table or a TikZ picture; no image file is copied into this package, the package declares no asset dependency and no image file is redistributed with it. Licence: version arXiv:2506.09609v2 of this article is distributed under the Creative Commons Attribution 4.0 International licence, as recorded in the arXiv licence metadata for this exact version and shown on the abstract page https://arxiv.org/abs/2506.09609v2. A full rights notice is delivered at licenses/arxiv-2506.09609-CC-BY-4.0.txt. Characters: every retained excerpt is pure ASCII, so the delivered engine, XeLaTeX with its default Latin Modern text font, reports no missing character.
\begin{enumerate}
    \item \label{condition:fractal-1} $\mathcal{A}_m \subset \varrho N^m \mathbb{Z}^2$ for $m \geq 0$;
    
    \item \label{condition:fractal-2} For $m \geq 0$, each vertex in $\mathcal{A}_m$ has $|\cdot|_\infty$-distance at least $\frac{1}{10} \varrho N^{m+1}$ away from each other.
\end{enumerate}

\begin{lemma}\label{lem:length-bound}
    There exists a constant $C = C(\varrho) \geq 100$ such that for all $N \geq 100 C$ and $k \geq 0$, all self-avoiding connected paths at level $k$ with $|\cdot|_\infty$-diameter $D \geq N^k$ have length at least $\frac{1}{C}(1 - \frac{C}{N})^kD$.
\end{lemma}

\begin{proof}
    We set $C = 10000 \varrho$ and prove the lemma by induction on $k$. The base case $k = 0$ holds trivially. For the inductive step, assume the statement holds for some $k \ge 0$. To prove the case $k+1$, consider any self-avoiding connected path $(x_1, x_2, \ldots, x_L)$ at level $(k+1)$ with diameter $D \geq N^{k+1}$. We divide it into sub-paths using the boxes from $\mathcal{S}_{k+1}$:
    $$ P_0, y_1, P_1, y_2, \ldots, y_M, P_M, $$
    where $y_1,\ldots,y_M \in \mathcal S_{k+1}$ and $P_0,\ldots,P_M$ are self-avoiding connected paths at level $k$ (with $P_0$ and $P_M$ possibly empty). Then,
    \begin{equation}\label{eq:lem5.13-1}
        {\rm diam}(P_0) + (2 \times 40\varrho + 20\varrho N^k) + {\rm diam}(P_1) + (2 \times 40\varrho + 20\varrho N^k) + \cdots +{\rm diam}(P_M) \geq D.
    \end{equation}
    We claim that for $N \geq 100 C$, 
    \begin{equation}\label{eq:lem5.13-2}
        \tfrac{C}{2N} ({\rm diam}(P_0) + {\rm diam}(P_1) + \cdots +  {\rm diam}(P_M)) \geq M (2 \times 40\varrho + 20\varrho N^k).
    \end{equation}
    We separate into two cases $M = 1$ and $M \geq 2$ (the case $M = 0$ is trivial). If $M \geq 2$, for each $1 \leq i \leq M-1$, since the $|\cdot|_\infty$-distance between the centers of $y_i$ and $y_{i+1}$ is at least $\frac{1}{10} \varrho N^{k+1}$ (Condition~\ref{condition:fractal-2}), we have ${\rm diam}(P_i) \geq \frac{1}{10} \varrho N^{k+1} - 20\varrho N^k - 2 \times 40 \varrho\geq \frac{1}{20} \varrho N^{k+1}$. This implies that
    $$
    \tfrac{C}{2N} ({\rm diam}(P_0) + {\rm diam}(P_1) + \cdots +  {\rm diam}(P_M)) \geq \tfrac{C}{2N} (M-1) \tfrac{1}{20} \varrho N^{k+1} \geq 100 M\varrho N^k \geq M(2 \times 40\varrho + 20\varrho N^k).
    $$
    If $M = 1$, we have $\min \{ {\rm diam}(P_0), {\rm diam}(P_M) \} \geq \frac{1}{3} N^{k+1}$, which leads to
    $$
    \tfrac{C}{2N} ({\rm diam}(P_0) + {\rm diam}(P_1) + \cdots +  {\rm diam}(P_M)) \geq \tfrac{C}{2N} \times \tfrac{1}{3} N^{k+1} \geq 2 \times 40\varrho + 20\varrho N^k.
    $$
    Combining both cases establishes Claim~\eqref{eq:lem5.13-2}.
    
    Let $A = \frac{1}{C} (1 - \frac{C}{N})^k$. For $1 \leq i \leq M-1$, recall that ${\rm diam}(P_i) \geq \frac{1}{20} \varrho N^{k+1} \geq N^k$. By the induction hypothesis, ${\rm len}(P_i) \geq A \cdot {\rm diam}(P_i)$ holds for all $1 \leq i \leq M - 1$. Similarly, we can also show that ${\rm len}(P_0) + A N^k \geq A \cdot {\rm diam}(P_0)$ and ${\rm len}(P_M) + A N^k \geq A \cdot {\rm diam}(P_M)$. Combining these estimates with~\eqref{eq:lem5.13-1} and~\eqref{eq:lem5.13-2} yields
    \begin{align*}
        AD &\leq A (1 + \tfrac{C}{2N}) ({\rm diam}(P_0) + {\rm diam}(P_1) + \cdots +  {\rm diam}(P_M)) \\
        &\leq (1 + \tfrac{C}{2N}) (2A N^k + {\rm len}(P_0) + {\rm len}(P_1) + \cdots +  {\rm len}(P_M)).
    \end{align*}
    Therefore, the total length $L$ satisfies
    \[ L \geq \sum_{i=0}^M {\rm len}(P_i) \geq (1 + \tfrac{C}{2N})^{-1} AD - 2A N^k. \]
    Since $C \geq 100$ and $N \geq 100C$, this yields $L \ge (1 - \frac{C}{N}) A D$ whenever $D \geq N^{k+1}$, thereby completing the induction.
\end{proof}

\end{document}
